% Folder 156-9, batch 07 — archivists' pages 121–139.
% Pass: Opus 5.5 (claude-opus-5-5), 2026-10-03 — first pass, unchecked against the pages by a human.
%
% Page 128 is a blank coloured sheet and is not transcribed. Page 129 opens a
% new run, paginated by him 1–11 (pp. 129–139), headed « Morphismes d'ateliers ».
\documentclass[11pt,a4paper]{article}
\input{../preamble/grothendieck}

\folder{156-9}
\batch{7}
\pages{121}{139}
\foldertitle{[Chapitre] IX et IX bis. [Ateliers] : notes manuscrites (05-15/07/1986).}
\dating{1986}
\watermark{Édition de démonstration}

\begin{document}

\page{121}
\note{p.~\uncertain{116} de l'auteur. La page commence au milieu d'une phrase : l'argument vient de la page précédente, hors de ce lot.}
de $A$ forment directement un \uncertain{At} \ill{}. (\uncertain{Ceci posé}) pour que le diagramme ($\mathcal{A}'$) soit commensurable i.e.\ $\bigcup A$ \add{\ill{}} \uncertain{commensurable}, il \ill{} et il suffit que $A$ le soit (i.e.\ qu'il soit « admissible », mais je préfère dire maintenant « commensurable »).

Nous pouvons maintenant poser

\textit{Axiome} \textit{Atmagsupp}~3 (englobant les deux précédents). Pour tout diagramme
\[
  F_0 \ll F_1 \ll \cdots \ll F_n ,
\]
le diagramme correspondant
\[
  A_0 \ll A_1 \ll \cdots \ll A_n
\]
dans $\operatorname{Figspatdisc}(\Sigma)$ \quad ($\Sigma = \Sigma_A$) est commensurable.

NB Nous pouvons, (pour une opération \uncertain{pseudo-$\Sigma$}) introduire le \ill{} atelier
\[
  \operatorname{Figélspatcomm}(\Sigma)
\]
des \uncertain{figures élémentaires} spatiales \emph{commensurables} de $\Sigma$, et \struck{des} \ill{} magiques \ill{}

\page{122}
\note{p.~117 de l'auteur.}
l'atelier $\operatorname{Figspatcomm}(\Sigma)$. Nous y prendrons $\trianglelefteq$, $|{\circ}|$ comme avant, mais une relation $\mathring{\ll}$ (ou je préfère écrire $\mathring{\ll}_{\mathrm{comm}}$, \ill{} \ill{} de \ill{}) plus forte.

Mais le sens (commensurable), i.e.\ que \struck{dépend} s'exprime peut-être \uncertain{à} l'aide des seules strates des figures $F$ et $G$, et que \ill{} sur les strates, il n'est pas clair que $\mathring{\ll}_{\mathrm{comm}}$ soit transitif ! (Déjà dans le cas trivial des figures élémentaires réduites à une seule « strate », via $g \leq f \leq e$
\[
  e = f \vee f' \text{ avec } f \mathrel{|{\circ}|} f', \qquad f = g \vee g' \text{ avec } g \mathrel{|{\circ}|} g',
\]
\note{la relation entre $f$ et $f'$, et entre $g$ et $g'$, est lue $|{\circ}|$ sans certitude.}
et donc \ill{} $g, g', f'$ mutuellement disjoints et \struck{$\{g \vee g', f'\}$} $\{g, g'\}$ ainsi que $\{g \vee g', f'\}$ commensurables, il ne \uncertain{s'ensuit} pas \ill{} \ill{} que $(g, g', f')$ l'\ill{}…) Donc ça ne semble pas marcher…

Ça \uncertain{marche} bien \uncertain{sûr} quand $\Sigma$ lui-même est une alg. de Boole, i.e.\ $\wedge$ est distributif \uncertain{par r.} à $\vee$…).

\page{123}
\note{p.~118 de l'auteur. Le premier alinéa est marqué d'un trait vertical dans la marge.}
Il faudrait voir si l'existence d'une « fonction support \uncertain{entière} fidèle » (cf \uncertain{VIII} p.~62) implique toujours la validité de \textit{Atmagsupp}~3, (comparer cf \uncertain{VIII} p.~78, cor.~1), et regarder en particulier le cas des \uncertain{magmas} fidèles.

Mais avant tout, il me faudrait \uncertain{aller} de l'avant, pour me rendre compte exactement comment vont intervenir les axiomes des supports, et s'ils \add{s'avèrent} \struck{\ill{}} aussi cruciaux que je le pressens.

\subsection*{Critique philosophique de cette approche}

\note{le titre est le début de la phrase ; il ne porte pas de soulignement distinct.}
\ill{} les supports, via les « opérations » \uncertain{vers} les « espaces ». Dans la définition même des opérations \uncertain{matérielles} \ill{}, \ill{}, on est forcé de « couper » des strates qui ne sont pas \uncertain{nécessairement} compatibles, (ni même « commensurables » \ill{}, et encore moins \uncertain{respectivement} disons, et encore moins combinatoirement). Or il semble que

\page{124}
\note{p.~119 de l'auteur.}
\uncertain{fin} de compte, je me suis intéressé aux supports et opérations sur les supports, que dans le cas des supports constructibles (cas « ens. constr. » — \uncertain{magmas} \uncertain{\ill{}} ($\operatorname{Supp}^{\circ} \mathcal{C}$…), et des \uncertain{familles} \uncertain{communes} de tels supports.

Il y a \struck{\ill{}} là une contradiction, qui est un \uncertain{visible} pour \uncertain{certains}, qui \ill{} \ill{} \ill{} \uncertain{maintenant}. D'ailleurs, \ill{} la \uncertain{dernière} \uncertain{présentation} $X \mathrel{|{\circ}|} Y$ implique \uncertain{que} la \uncertain{connaissance} de \uncertain{strates} $X$, $Y$ qui le \uncertain{composent} \uncertain{quelles} \uncertain{sont} « \uncertain{ind.} disj. ») \ill{} \ill{} \ill{} \uncertain{peuvent} ne pas être compatibles, ni même commensurables.

Pourtant, il semble indispensable \uncertain{d'avoir} une description « \uncertain{intrinsèque} » des \uncertain{supports} \uncertain{subordonnés} à une figure, (\ill{} \ill{} \ill{} \uncertain{objets} de \uncertain{la} figure \ill{} l'occurrence \ill{} $\Sigma_{\mathcal{A}}$) indépendamment (\ill{} \ill{} \ill{}) \ill{} de la dite figure. On \ill{} \uncertain{plus}

\page{125}
\note{p.~120 de l'auteur.}
\uncertain{aussi} \struck{\ill{}} devrait-il trouver, pour une fois ? Il faudrait que j'y sois très attentif, chaque fois que je \uncertain{ferai} appel à des considérations de support. Je voudrais p.~ex.\ donner aux supports \uncertain{attachés} à une figure, et leurs relations aux supports \uncertain{attachés} à une subdivision, une \uncertain{définition} « \uncertain{intrinsèque} » qui ne soit affectée par le fait que les deux subdivisions ne soient \uncertain{pas} « commensurables » — ce « \uncertain{sens} » un peu vague peut-il se préciser ?

On pourrait bien sûr \uncertain{songer} \uncertain{à} \uncertain{introduire} plutôt \uncertain{que} \ill{} \ill{} \ill{} \ill{} $\Sigma_{\mathrm{ens}}$ de $\Sigma = \Sigma_{\mathcal{A}}$, formé \ill{} des supports ens.\ constructibles, \ill{} des opérations $\mathcal{S}$ (partout définies) et \uncertain{des} \ill{} et \ill{} \uncertain{partiellement} définies seulement). \struck{\ill{}} \uncertain{Cette} \ill{} \uncertain{explicite}-il une construction indépendante de la considération

\page{126}
\note{p.~121 de l'auteur.}
\uncertain{notion} préalable de $\Sigma$, le \uncertain{vrai} \uncertain{chemin} ? Par exemple, les supports \add{constr.} \uncertain{fermés} (\struck{\uncertain{exclusif}} \add{\uncertain{relatif}} à un \uncertain{magma} \ill{} $\mathcal{F}$) pourraient-ils s'obtenir, en prenant l'ensemble \uncertain{ordonné} déduit de $(\mathcal{F}, \ll)$ en « inversant » les « flèches » $\ll$ (\ill{} les \ill{} subdivisions) ? \ill{} \ill{} \ill{} pas sûr que la cat.\ des fractions obtenue soit bien un ens.\ ordonné, (\ill{} s'il l'était) et encore moins qu'il \struck{\ill{} \ill{}} s'identifie à un sous-ens.\ ordonné de $\Sigma$. (Ce \ill{} \ill{} nécessairement un \ill{} \ill{} si $\Sigma$ devait être \uncertain{lié} aux \uncertain{profils} des formes…).

Je dirais quand-même, pour me \uncertain{rassurer}, que dans les cas où $\Sigma$ est booléen (ateliers \uncertain{pseudo-booléens}, et peut-être les ateliers fidèles ?, en tout cas les ateliers \uncertain{magmatiques}), il n'y a \uncertain{vraiment} rien à lui reprocher. Ce \uncertain{incluant} \uncertain{avec} \ill{} déjà \uncertain{tous} les cas qui \uncertain{importent}.

\page{127}
\note{p.~122 de l'auteur.}
considérer \struck{\ill{}} soit pas « \uncertain{classiques} ». Ce sont des « ordres », pseudo-simpliciaux \uncertain{topologiques} et \uncertain{tutti quanti}, associés à l'intuition du « lieu » remplaçant le « point », qui causent tous les \uncertain{ennuis} pour l'instant, et les seuls ennuis.

Mais même dans le cas des ateliers \uncertain{pseudo-booléens}, il serait intéressant de voir si $\Sigma_{\mathrm{ens}}$ n'admet pas une description combinatoire plus proche directe, dans un esprit plus proche d'une « géométrie des formes » — dans laquelle la \uncertain{spatialité} a des \uncertain{fortes} relents \ill{} de topologie générale !\note{cet alinéa est marqué d'un trait vertical dans la marge.}

Prochaines choses : regarder pour « boucler » l'aspect provisoirement ensembliste — il serait temps !
\begin{itemize}
  \item[a)] Subdivisions — enfin ! \add{morphismes}\note{« morphismes » est écrit au-dessus de « Prochaines choses », souligné, et relié par un trait à « Subdivisions ».}
  \item[b)] Ateliers modérés
  \item[c)] Axiomes de saturation. Saturation d'un atelier, ou d'un atelier magmatique…
\end{itemize}

\marginal{Avant, peut-être, \uncertain{a)} \ill{} sur les ateliers, \uncertain{b)} \ill{} relatifs \uncertain{aux} supports}\note{note écrite verticalement dans la marge gauche, en face de la liste ; ses lettres a), b) renvoient aux points de la liste.}

\page{129}
\note{p.~1 de l'auteur, qui inscrit en tête « cf IX bis » ; une nouvelle pagination commence ici. Dans le coin supérieur gauche, en diagonale : « cf \uncertain{VIII} ».}
\section*{Morphismes d'ateliers}

\marginal{\ill{}}\note{quelques mots en diagonale dans la marge gauche, en face du premier alinéa, illisibles.}

\struck{\ill{}} Je suis intéressé par trois types de situations, que je voudrais pouvoir exprimer par une notion de morphismes

1°) Cas d'un atelier \add{$\mathcal{A}$}, et d'un sous-atelier $\mathcal{A}'$, i.e.\ (cf \uncertain{X} p.~27) une partie $\mathcal{A}'$ satisfaisant la condition

\textit{Sous at} Si $X \mathrel{\mathring{\ll}} Y' \trianglelefteq Y$, avec $X, Y \in \mathcal{A}'$, alors $Y' \in \mathcal{A}'$.

Cela implique que, muni des relations induites par $\mathcal{A}$ $\trianglelefteq$, $\mathring{\ll}$, $|{\circ}|$, $\mathcal{A}'$ est un atelier, et que de plus, la relation $\ll_{\mathcal{A}'}$ est induite par $\ll_{\mathcal{A}}$ (ce qui est assuré précisément par Sous at). Je rappelle les exemples principaux

\textit{Ex 1} Soit $\mathcal{F} \subset \mathcal{A}$ une préfigure de $\mathcal{A}$, i.e.\ une partie telle que $\overline{\mathcal{F}}^{(\trianglelefteq)}$ soit une figure, i.e.\ telle que $X, Y \in \mathcal{F} \Rightarrow X \between Y$.\note{le signe noté ici $\between$ (une croix barrée) est rendu ainsi faute de mieux, dans tout le lot ; sa forme exacte reste douteuse.} Alors
\[
  \operatorname{Omb}^{\circ}(\mathcal{F}) = \{ X \in \mathcal{A} \mid X \mathrel{\mathring{\ll}} \mathcal{F} \}
\]
est un sous-atelier. (loc.\ cit.)

\textit{Ex 1 bis} $\mathcal{F}$ \struck{\ill{}} \add{\uncertain{comme}} figure de $\mathcal{A}$, alors $\mathcal{F}$ lui-même est un sous-atelier, (qui est \uncertain{même} \emph{magmatique} pour la structure induite). (En effet, si $X \mathrel{\mathring{\ll}} Y' \trianglelefteq Y$, avec $X, Y \in \mathcal{F} \subset \mathcal{F}$, \struck{\ill{}} comme $Y' \in \overline{\mathcal{F}}$, on a $X = Y'$, donc $Y' \in \mathcal{F}$.)

\page{130}
\note{p.~2 de l'auteur.}
\textit{Ex 2} Soit $S \subset \mathcal{A}$, considérons « l'atelier résiduel strict » formé des $X \in \mathcal{A}$ tels que $X \mathrel{\|} S$, il est \struck{\ill{}} \add{\uncertain{fermé}} pour $\ll$ et a fortiori pour $\trianglelefteq$, donc c'est un sous-atelier.

\textit{Exemple 2 bis} Considérons par contre le sous-ensemble
\[
  \operatorname{Supp}^{\circ}(S) = \{ X \in \mathcal{A} \mid X \mathrel{|{\circ}|} S \}.
\]
Ce n'est pas en général un sous-atelier (loc.\ cit.\ p.~28). Mais dans le cas où $S = \mathcal{F}$ est une préfigure, considérons l'ensemble \struck{des}
\[
  \{ X \in \mathcal{A} \mid X \mathrel{|{\circ}|} \mathcal{F} \text{ et } X \between \mathcal{F} \}
\]
(atelier résiduel au sens large). C'est un sous-atelier. En effet, soit
\[
  X \mathrel{\mathring{\ll}} Y' \trianglelefteq Y \quad \text{avec } X, Y \in \mathcal{A}',
\]
prouvons $Y' \in \mathcal{A}'$, i.e.
\[
  \forall Z \in \mathcal{F}, \quad Y' \between Z, \text{ et } Y' \mathrel{|{\circ}|} Z \text{ i.e. } Y' \neq Z .
\]
La relation \struck{$Y' \between Z$ claire car $Y \between Z$}, \add{$Y' \trianglelefteq Y$…} \uncertain{Si} \uncertain{on} avait $Y' = Z$, on aurait $X \mathrel{\mathring{\ll}} Z$, or $X \in \mathcal{A}'$ implique que $X \mathrel{|{\circ}|} Z$, donc on aurait $X \mathrel{|{\circ}|} X$, absurde.\note{le passage rayé est souligné en pointillé, ce qui pourrait indiquer qu'il est rétabli ; « $Y' \trianglelefteq Y$… » est écrit au-dessus.}

Dans un \uncertain{prochain} cas, il faudrait vérifier si $\Sigma_{\mathcal{A}'} \hookrightarrow \Sigma_{\mathcal{A}}$ ?\note{alinéa marqué d'un double trait dans la marge.}


\page{131}
\note{p.~3 de l'auteur.}
2°) Cas d'une application ensembliste \struck{de} $f : E' \to E$, ou d'espaces modérés, \add{\ill{}} et de l'application
\[
  f^{*} : \operatorname{Figures}(E) \longrightarrow \operatorname{Figures}(E')
\]
\uncertain{correspondante}. Si $f$ n'est pas surjective, cette application n'applique pas $\operatorname{Figelann}(E) = \mathcal{A}$ dans $\operatorname{Figelann}(E') = \mathcal{A}'$, mais dans \struck{$\operatorname{Figelann} \cup \{\emptyset\}$} $\mathcal{A}' \cup \{\emptyset\}$ \quad ($\leftarrow$ figure vide).


\marginal{\ill{}}\note{note marginale en diagonale dans la marge gauche, en face des lignes sur $f_{*}$ ; quelques mots seulement, dont \uncertain{d'application}, sont distinguables.}

Cette application est compatible avec $\trianglelefteq$, $\mathring{\ll}$, $|{\circ}|$, mais en général elle n'est pas injective… Elle a tendance à passer aux supports, et donc à commuter à tout ($\mathrm{Sup}$, $\mathrm{Inf}$, $\complement$).

3°) Cas d'une application ensembliste $g : E \to E'$, injective, et application
\[
  g_{*} : \text{\struck{$\operatorname{Figel}$}}\ \mathcal{A} \longrightarrow \mathcal{A}', \quad \text{\uncertain{alors} injective.}
\]
Commute avec : $\trianglelefteq$, $\mathring{\ll}$, $|{\circ}|$, mais en passant aux supports, commute à $\operatorname{Sup}_{i \in I}$ ($I$ quelc.), $\operatorname{Inf}_{i \in I}$ ($I \neq \emptyset$), $S \wedge \complement T$, mais pas à $\complement$ si $g$ pas bijective.


\page{132}
\note{p.~4 de l'auteur.}
Cas commun : 2° et 3° : correspondance

\begin{tikzcd}
\Gamma \arrow[r, hook, "f'"] \arrow[d, "f"'] & E' \\
E &
\end{tikzcd}

telle que $\Gamma \to E'$ soit injectif. Dans ce cas \ill{} 2° et 3° (en fait, \ill{} \ill{}) \ill{}. \ill{}

\marginal{Le test \ill{} \ill{} dans le cas 2°, est $\hat{f}_{*}(1)=1$. Le test pour le cas 3° est quoi ?? \ill{} Pour l'injectivité \ill{} (\ill{} \ill{}), \ill{} Cela fait que $f$ soit injectif, \ill{} \ill{} \ill{} $\mathcal{A}$ \ill{} \ill{} $|{\circ}|$ (\ill{} \ill{} commutation \ill{}) \ill{} le fait que $\beta_{*}$ \ill{} l'application \ill{} l'injectivité de \ill{} \ill{} $\mathrm{Sup}$, $\complement T$ … \ill{} \ill{} \uncertain{injective}.}\note{longue note en diagonale qui envahit la moitié gauche de la page ; lue en tournant la feuille, elle reste en grande partie illisible.}

\ill{} $\mathcal{A} \to \mathcal{A}' \cup \{\emptyset\}$.

Il \uncertain{semble} qu'il n'y a pas d'application directe de figures, par une application $g : E \to E'$ qui ne soit injective.

\textit{Dans le cas d'un morphisme d'ateliers quelconque}, il faut \uncertain{voir si} \ill{} se réduit toujours \uncertain{aux} « correspondances » sur les opérations \uncertain{induites} — \ill{} $\alpha_{*}$, $\beta_{*}$.

\page{133}
\note{p.~5 de l'auteur.}
\textit{Définition} (provisoire). Soient $\mathcal{A}$, $\mathcal{A}'$ deux ateliers. Un \struck{ho} homomorphisme (algébrique, sans connotation géométrique fixée, pour le moment) de $\mathcal{A}$ dans $\mathcal{A}'$, est une application \struck{de $\mathcal{A}$}
\[
  f : \mathcal{A} \longrightarrow \hat{\mathcal{A}}' := \mathcal{A}' \amalg \{\emptyset\}
  \qquad \bigl(\hookrightarrow \mathcal{F}(\mathcal{A}')\bigr)
\]
satisfaisant les trois conditions, pour $X, Y \in \mathcal{A}$ :
\begin{itemize}
  \item[a)] $X \trianglelefteq Y \Longrightarrow f(X) \trianglelefteq f(Y)$
  \item[b)] $X \mathrel{\mathring{\ll}} Y \Longrightarrow f(X) \mathrel{\mathring{\ll}} f(Y)$
  \item[c)] $X \mathrel{|{\circ}|} Y \Longrightarrow f(X) \mathrel{|{\circ}|} f(Y)$.
\end{itemize}
\note{les lettres a), b), c) sont écrites sur des numéros i), ii), iii).}

\marginal{Mais la condition \ill{} cf \uncertain{VIII} \ill{} \ill{}}

Pour donner un sens à cette définition, il faut donner un sens aux relations
\[
  X' \trianglelefteq Y', \quad X' \mathrel{\mathring{\ll}} Y', \quad X' \mathrel{|{\circ}|} Y'
  \quad \text{pour } X', Y' \in \mathcal{A}' \amalg \{\emptyset\},
\]
donc \uncertain{définir} quand l'un ou l'autre des deux termes est \struck{\ill{}} égal à $\emptyset$. Bonne convention ici pour $\forall X' \in \mathcal{A}'$ :

a) $X' \trianglelefteq \emptyset$, \struck{\ill{}} ($\emptyset \trianglelefteq X'$ \struck{\ill{}} toujours vrai) (NB $\trianglelefteq$ sur $\hat{\mathcal{A}}'$ \struck{\ill{}} n'est pas une relation d'ordre !)

b) $\emptyset \mathrel{\mathring{\ll}} X'$, ($X' \mathrel{\mathring{\ll}} \emptyset \iff X' = \emptyset$)

NB $\mathring{\ll}$ sur $\hat{\mathcal{A}}'$ est une rel. d'ordre admettant \uncertain{$\emptyset$} comme plus petit él.

\page{134}
\note{p.~6 de l'auteur.}
c) $X' \mathrel{|{\circ}|} \emptyset$ toujours satisfait.

L'introduction de $\hat{\mathcal{A}}'$, et \uncertain{des} conventions précédentes, sont inspirées par l'exemple 2°) précédent.

\textit{Noyau} de $f$ :
\[
  \operatorname{Ker} f \overset{\text{déf}}{=} \{ X \in \mathcal{A} \mid f(X) = \emptyset \}.
\]
Alors b) implique que $\operatorname{Ker} f \subset \mathcal{A}$ est fermé par $\mathring{\ll}$. Du coup on s'attend que ce soit des $\Sigma(\mathcal{A})$, mais on doit \uncertain{sans doute} pour cela imposer des conditions supplémentaires sur $f$ (en tous cas c'est OK dans l'exemple 2° \ill{} \ill{}).

Partant d'une partie \struck{$N$} $\subset \mathcal{A}$, fermée pour $\mathring{\ll}$, est-ce \uncertain{le} noyau d'un homomorphisme d'atelier, et \uncertain{ce dernier} peut-il se choisir canoniquement ? Plus précisément, peut-on trouver un atelier $\mathcal{A}/N$ (opération d'\emph{excision})

\page{135}
\note{p.~7 de l'auteur.}
et un homom.
\[
  \mathcal{A} \xrightarrow{\;p\;} \mathcal{A}/N
\]
ayant le noyau $N$ (ou $\supset N$ ?), qui soit \emph{universel}, au sens que tout homom. d'ateliers
\[
  \mathcal{A} \xrightarrow{\;f\;} \mathcal{A}'
\]
ayant un noyau $\supset N$, se factorise par $p$ de façon unique. (NB La notion de \emph{composition} est claire, en interprétant les homom. d'ateliers comme des applications ensemblistes
\[
  \mathcal{A} \amalg \{\emptyset\} \xrightarrow{\;\hat{f}\;} \mathcal{A}' \amalg \{\emptyset\}
\]
transformant $\{\emptyset\}$ en $\{\emptyset\}$, et \struck{\ill{}} \struck{pour} dans la construction \ill{} $\mathcal{A}$ \ill{} les propriétés \ill{}. \ill{} \ill{} \ill{} dire que $\hat{f}$ est compatible à $\trianglelefteq$, $\mathring{\ll}$, $|{\circ}|$ sur $\hat{\mathcal{A}}$ !).

\marginal{\ill{}}\note{quelques mots en diagonale dans la marge gauche, illisibles.}

L'exemple type est celui-ci :
\[
  \mathcal{A} = \operatorname{Figelann}(E), \qquad N = \{ X \in \mathcal{A} \mid |X| \subset E \setminus E' \},
\]
\[
  \mathcal{A}/N = \operatorname{Figelann}(E'),
\]
mais \ill{} \ill{} \ill{} \ill{} \uncertain{comme} \uncertain{quotient de} $\mathcal{A} \setminus N$.\note{fin de page très chargée : sous « $\operatorname{Figelann}(E)$ » il ajoute \uncertain{dans ce cas}, et une ligne d'interligne, en partie rayée, court jusqu'à « $\mathcal{A}\setminus N$ ». Le nom $\operatorname{Figelann}$ est lu tel qu'il est écrit, sans certitude sur sa graphie.}

\page{136}
\note{p.~8 de l'auteur.}
À supposer qu'une telle construction existe dans un cas donné \ill{} pas \uncertain{général}, elle s'annonce \ill{} \uncertain{délicate}. Pour fixer les idées, il serait bon \struck{\ill{}} \ill{} \ill{} regarder \uncertain{d'abord} les cas suivants

1°) Cas d'un \struck{\ill{}} atelier ensembliste

\struck{\ill{}} Déjà ce n'a pas l'air évident, surtout \ill{} \ill{} on \ill{} suppose pas que l'atelier soit modéré, \struck{\ill{}} \uncertain{même} pour des $N$ très particuliers comme p. ex. $\operatorname{Omb}^{\circ} F$, $F$ une figure, \uncertain{ou} $\operatorname{Supp} F$.

2°) Cas d'un \uncertain{ensemble ordonné} $A$, $N$ partie quelconque de $A$ (NB $x \mathrel{\mathring{\ll}} y \Rightarrow x = y$ !). \struck{L}Là, on doit trouver l'\uncertain{ensemble ordonné} $A \setminus N$ simplement. \ill{} \ill{} \ill{} cette fois !

Pour la relation avec le formalisme des \uncertain{supports}, s'inspirer éventuellement de cf \uncertain{VIII} p.~70 ff. Là je n'ai pas le \ill{} pour, et laisse tomber pour le moment.

\page{137}
\note{p.~9 de l'auteur.}
\subsection*{Subdivisions}

Soient $\mathcal{A}$ un \struck{\ill{}} \add{atelier}, $\mathcal{F}$ l'ensemble des $\mathcal{A}$-figures, et
\[
  G \ll F
\]
dans $\mathcal{F}$. On a donc une application \emph{croissante}
\[
  \varphi : \widetilde{G} \longrightarrow \widetilde{F}
\]
\note{sous $\widetilde{G}$ et $\widetilde{F}$ il écrit, par un signe d'égalité vertical, $G$ et $F$ munis d'un signe diacritique qu'on ne lit pas sûrement.}
telles que pour $Y \in \widetilde{G}$, $X = \varphi(Y)$, on ait
\[
  Y \mathrel{\mathring{\ll}} X .
\]
Pour dire que $G$ est une \struck{\ill{}} \emph{subdivision} de $F$, nous exigeons \struck{d'abord} que

(i) $\forall X \in \widetilde{F}$, \struck{posant} considérons $\varphi^{-1}(X) \subset \widetilde{G}$, on a $\varphi^{-1}(X) \mathrel{\mathring{\ll}} X$, et $\operatorname{Supp}^{\circ} \varphi^{-1}(X) \subset X^{\circ}$, \ill{} \struck{on a} \add{donc} \struck{$\operatorname{Supp}^{\circ} \varphi^{-1}(X)$}
\[
  \operatorname{Supp}^{\circ} \varphi^{-1}(X) = X^{\circ}
\]

\marginal{Il \uncertain{valait mieux} ici \ill{} \struck{\ill{}} \ill{} \ill{} \ill{}}\note{note en diagonale dans la marge gauche, en partie rayée.}

\uncertain{Atmag supp} ?, on voit facilement que ceci équivaut à chacune des deux conditions

(ii) $\operatorname{Supp} G = \operatorname{Supp} F$ \quad ($= \operatorname{Supp}^{\circ} G$) \quad ($= \operatorname{Supp}^{\circ} F$)

ou encore

(iii) Il n'existe pas $Z \in \mathcal{A}$ avec $Z \ll F$, et $Z \mathrel{|{\circ}|} G$.

\page{138}
\note{p.~10 de l'auteur.}
Quand $\mathcal{A}$ est divisible, \struck{\ill{}} $Z \in \mathcal{L}$, et on trouve la forme

\struck{(iii)} \add{(iv)} Il n'existe pas $x \in \mathcal{L}$, avec $x \in \operatorname{Omb}(F)$ (i.e.\ $x \ll F$), et $x \mathrel{|{\circ}|} G$.

En fait, les implications
\[
  \text{(i)} \Rightarrow \text{(ii)} \Rightarrow \text{(iii)} \Rightarrow \text{(iv)}
\]
sont triviales, et (iv) $\Rightarrow$ (i) (\ill{} \ill{} divisibilité) aussi ! (iii) $\Rightarrow$ (i), qu'en \ill{} \uncertain{Atmag supp} ?\note{dans cette chaîne, plusieurs numéros sont récrits sur d'autres, rayés.}

[Signalons \ill{} une autre condition équivalente, moyennant la \emph{divisibilité} :

\struck{(iv \ill{})} \add{(v)} $G$ est maximal parmi les raffinements de $F$, admis par \struck{\ill{}} $\mathcal{C}$.

\struck{Mais} \add{Quand} ces conditions \struck{\ill{}} \add{\ill{}} \ill{} \struck{satisfaites}, je dirai \struck{\ill{}} que $G$ est une \emph{subdivision de} $F$. \struck{(comme précédemment)}

\marginal{\ill{}}\note{note en diagonale dans la marge gauche, illisible ; on y distingue un « NB ».}

\struck{Mais je rajouterai la condition}

\struck{(ii) La top. de $\widetilde{F}$ (\ill{} d'ordre) est \uncertain{quotient} de celle de $\widetilde{G}$, \ill{} en \uncertain{d'autres termes}, si $X, X' \in \widetilde{F}$, pour qu'on ait $X \leq X'$, il faut (et il suffit) qu'il existe $Y \in \varphi^{-1}(X)$, $Y' \in \varphi^{-1}(X')$ avec $Y \leq Y'$.}

\struck{Cela signifie aussi, dans le cas $\mathcal{A} = \operatorname{Figelann}(E)$}…\note{ce dernier passage est barré de traits diagonaux ; la rature se poursuit sur la page suivante.}

\page{139}
\note{p.~11 de l'auteur.}
\struck{disons, que $F$ se déduit de $G$ par « regroupement de \uncertain{strates} » $Y_j^{\circ}$ (pour former les $X_i^{\circ}$) sans plus -- et que la relation d'ordre sur $I$, \ill{} ces « strates fermées », se déduit simplement de celle sur $J$ par passage au quotient.}\note{page barrée de deux traits diagonaux, comme la fin de la précédente ; le reste de la feuille est vide. C'est la dernière page du dossier : le passage rayé s'arrête là.}

\end{document}
